\section{Simple Tables}\label{sec:simple-tables}

Two table formats are available in TAOS. The first is called the \emph{simple}
format because it is easy to create, but it has limited capability. The second is
called the \emph{full} format because it provides the full capabilities of the
software, but it is somewhat more difficult to create.

Simple tables consist of a single set of tabulated numbers. The set defines a function
of one or more state variables, or independent variables. The simplest table is a
function of one variable. For example, a table of thrust values as a function of time
is

\begin{lstlisting}[style=taosmanual]
(motor-1)
  table thrust(time)
    time   = 0.0, 0.5, 2.0, 4.0, 4.5, 5.0
    thrust = 0.0, 3500., 3100., 2950., 600., 0.0
\end{lstlisting}

The table begins with a name in parentheses followed by the keyword
\texttt{table}. This is followed by the table type, \tabletype{thrust}, and a list of
independent variables---in this case, \statevar{time}. Tabulated values for the
independent variable and the dependent variable are given next. For each value of
time, a corresponding thrust value is given.

\taossubhead{Table Identification Name}

Simple tables always begin with a name enclosed in parentheses. This name is the
\emph{table identification name}, and it is used in the problem file to reference the
table. The name must be unique among all input table files. It cannot contain spaces
or other delimiters; however, nondelimiter special characters such as underscores,
periods, or dashes are encouraged to improve readability.

\taossubhead{Table Function Definition}

The table identification name is followed by the keyword \texttt{table} and the table
type from \cref{tab:table-types}. A list of as many as five state variables, or
independent variables, follows. These variables must be separated by commas and
enclosed in parentheses so that the input resembles a function definition. For
example,

\begin{lstlisting}[style=taosmanual]
table ca(mach,alt,alpha)
\end{lstlisting}

says that $C_A$ is a function of Mach number, altitude, and angle of attack.

\taossubhead{Extrapolation Option}

The function definition is followed by an extrapolation keyword, either
\texttt{extrap} or \texttt{no-extrap}. This keyword is optional and controls how
TAOS handles independent-variable values outside the tabulated range, as shown in
\cref{fig:table-extrapolation-option}.

\begin{figure}[htbp]
  \centering
  \begin{tikzpicture}[x=0.72cm,y=0.72cm,>=Latex,font=\scriptsize]
  \newcommand{\gridbox}[1]{%
    \begin{scope}[shift={#1}]
      \draw[line width=0.75pt] (0,0) rectangle (5,5);
      \foreach \x in {1,2,3,4} {\draw[densely dotted] (\x,0)--(\x,5);}
      \foreach \y in {1,2,3,4} {\draw[densely dotted] (0,\y)--(5,\y);}
      \node[below,font=\small] at (2.5,0) {$x$};
      \node[left,font=\small] at (0,2.5) {$f(x,z)$};
      \node[left] at (0.12,2.15) {$z_1$};
      \node[left] at (0.12,1.02) {$z_2$};
    \end{scope}%
  }
  \gridbox{(0,0)}
  \gridbox{(8.0,0)}
  \node[font=\small] at (2.5,5.55) {Extrapolation};
  \node[font=\small] at (10.5,5.55) {No Extrapolation};

  \begin{scope}
    \draw[line width=0.85pt] (0.45,2.0) .. controls (2.0,1.95) and (3.0,2.35) .. (4.25,3.65);
    \draw[line width=0.85pt] (0.45,0.95) .. controls (1.8,1.0) and (3.0,1.4) .. (4.25,2.65);
    \foreach \p in {(0.45,2.0),(1.45,2.05),(2.5,2.35),(3.45,2.85),(4.1,3.45),
                      (0.45,0.95),(1.55,1.15),(2.5,1.55),(3.45,2.05),(4.1,2.45)}
      {\draw[fill=white] \p rectangle ++(0.20,0.20);}
    \draw[dashed,line width=0.8pt] (4.15,3.55)--(4.75,4.20);
    \draw[dashed,line width=0.8pt] (4.15,2.55)--(4.75,3.20);
    \draw[->] (4.60,4.05)--(4.60,1.15);
    \draw[fill=black] (4.45,1.1) rectangle ++(0.22,0.22);
    \draw[fill=white] (4.45,3.75) rectangle ++(0.22,0.22);
  \end{scope}

  \begin{scope}[shift={(8.0,0)}]
    \draw[line width=0.85pt] (0.45,2.0) .. controls (2.0,1.95) and (3.0,2.35) .. (4.25,3.65);
    \draw[line width=0.85pt] (0.45,0.95) .. controls (1.8,1.0) and (3.0,1.4) .. (4.25,2.65);
    \foreach \p in {(0.45,2.0),(1.45,2.05),(2.5,2.35),(3.45,2.85),
                      (0.45,0.95),(1.55,1.15),(2.5,1.55),(3.45,2.05)}
      {\draw[fill=white] \p rectangle ++(0.20,0.20);}
    \draw[dashed,line width=0.8pt] (4.0,3.25)--(4.6,3.95);
    \draw[dashed,line width=0.8pt] (4.0,2.25)--(4.6,2.95);
    \draw[fill=black] (3.85,2.2) rectangle ++(0.22,0.22);
    \draw[fill=white] (4.45,0.95) rectangle ++(0.22,0.22);
  \end{scope}
\end{tikzpicture}

  \caption{Table extrapolation option.}
  \label{fig:table-extrapolation-option}
\end{figure}

All interpolation and extrapolation is linear, which should be considered when
entering values and setting the extrapolation option. The default is to allow
extrapolation.

\taossubhead{Table Parameters}

The extrapolation option may be followed by a variable name and value called the
\emph{table parameter}. This parameter provides additional information associated
with the table.

For aerodynamic-coefficient tables, the parameter is the reference area,
\statevar{sref}. For thrust and mass-flow tables, a parameter called \texttt{units}
gives the units used in the table. Allowable units are listed in
\cref{tab:thrust-mass-flow-units}. Other table types do not use a table parameter.

\begin{table}[htbp]
  \centering
  \caption{Units for thrust and mass-flow tables.}
  \label{tab:thrust-mass-flow-units}
  \begin{tabular}{@{}llll@{}}
    \toprule
    Thrust units & Mass-flow units & Mass-flow units & Mass-flow units \\
    \midrule
    \texttt{lb} & \texttt{lb/sec} & \texttt{lb/min} & \texttt{lb/hr} \\
    \texttt{n} (Newtons) & \texttt{slugs/sec} & \texttt{slugs/min} & \texttt{slugs/hr} \\
    \texttt{kn} (kilonewtons) & \texttt{g/sec} & \texttt{g/min} & \texttt{g/hr} \\
    & \texttt{kg/sec} & \texttt{kg/min} & \texttt{kg/hr} \\
    \bottomrule
  \end{tabular}
\end{table}

\taossubhead{Examples}

A complete table-definition line for an axial-force-coefficient table is

\begin{lstlisting}[style=taosmanual]
table ca(mach,alpha) extrap sref=2.162
\end{lstlisting}

This defines a table that is a function of Mach number and angle of attack, permits
extrapolation, and uses a reference area of 2.162~ft\textsuperscript{2}. The units of
\statevar{sref} default to square feet, but they can be changed with the
\problemblock{*units/fmt} data block in the problem file
(\cref{sec:block-units-fmt}). Another example is

\begin{lstlisting}[style=taosmanual]
table mdot(time,motor) no-extrap units=lb/hr
\end{lstlisting}

which defines a mass-flow-rate table that is a function of time and a user-defined
variable called \texttt{motor}, does not allow extrapolation, and has units of
pounds per hour.

\taossubhead{Independent Variable Values}

The table-definition information is designed to appear on one line. It is followed by
values for the independent variables. The variable name is given first, followed by a
list of values, for example,

\begin{lstlisting}[style=taosmanual]
mach = 0.2, 0.4, 0.6, 0.8, 0.9
alpha = 0, 2, 4
\end{lstlisting}

Any number of values can be input; the number is limited only by the computer's
available memory. Values for independent variables must be strictly increasing or
strictly decreasing, and duplicate values are not allowed.

Values can be entered with or without decimal points, and all are treated as real
numbers. Scientific notation is also allowed:

\begin{lstlisting}[style=taosmanual]
time = 0., 1.0e1, 2.0e1, 6.0e1, 7.0e1, 9.0e1,
       2e02, 4e02, 6e02
\end{lstlisting}

Spaces or commas are usually used to separate values because they make the
separation clear. Values need not all appear on one line; they can be continued as in
the preceding example.

When a table is a function of more than one independent variable, the variables and
their values must be given in the same order as in the function definition. For
example,

\begin{lstlisting}[style=taosmanual]
table cx(alt,mach) extrap sref=5.30 # CORRECT:
  alt = 0, 10000                    # altitudes first,
  mach = 5, 10, 15, 20             # then Mach numbers
\end{lstlisting}

is correct because altitude precedes Mach number in both the function definition and
the value lists. An error is generated if the Mach-number values are given first:

\begin{lstlisting}[style=taosmanual]
table cx(alt,mach) extrap sref=5.30 # INCORRECT: altitude
  mach = 5, 10, 15, 20             # is first in definition,
  alt = 0, 10000                    # but Mach values are
                                    # given first
\end{lstlisting}

\taossubhead{Dependent Variable Values}

After all independent-variable values have been given, the dependent values are
listed. For tables with a single independent variable, there is a one-to-one
correspondence between the independent and dependent values. The following
complete thrust table illustrates this:

\begin{lstlisting}[style=taosmanual]
(example-thrust)
  table thrust(time) units=lb no-extrap
    time   = 0.0, 0.1, 1.0, 2.0, 3.0, 3.4, 3.5
    thrust = 0, 5100, 4950, 4900, 4840, 4700, 0
\end{lstlisting}

At \statevar{time} equal to \SI{2.0}{\second}, thrust is \SI{4900}{\poundforce}. The
same number of values must be given for both \statevar{time} and
\statevar{thrust}.

For tables with more than one independent variable, the dependent values must be
given in the correct order. For example,

\begin{lstlisting}[style=taosmanual]
(example-cx)
  table cx(alt,mach) extrap sref=5.30
    alt = 0, 10000
    mach = 5, 10, 15, 20
    cx = 0.030, 0.027, 0.025, 0.024 # Values for alt=0
         0.035, 0.031, 0.029, 0.028 # Values for alt=10000
\end{lstlisting}

Values for each Mach number at the first altitude are given first, followed by values
for the second altitude. A dependent value is required for every combination of
independent values. This example has two altitudes and four Mach numbers, so eight
$C_X$ values are required.

The following table is a function of three independent variables:

\begin{lstlisting}[style=taosmanual]
(example-ca)
  table ca(alt,mach,alpha) no-extrap sref=10.0
    alt = 0, 5000
    mach = 4, 6, 8
    alpha = 0, 2, 4, 6
    ca = 0.020, 0.021, 0.022, 0.024 # mach=4, alt=0
         0.018, 0.019, 0.020, 0.022 # mach=6, alt=0
         0.017, 0.018, 0.019, 0.021 # mach=8, alt=0

         0.021, 0.022, 0.023, 0.025 # mach=4, alt=5000
         0.020, 0.021, 0.022, 0.024 # mach=6, alt=5000
         0.019, 0.020, 0.021, 0.023 # mach=8, alt=5000
\end{lstlisting}

There are two altitudes, three Mach numbers, and four angles of attack, so there must
be $2\cdot3\cdot4=24$ values for $C_A$. Values for altitude zero are given first.
Within this set, values for Mach 4 are given first, followed by values for Mach 6, and
so on. This pattern continues for all values of Mach number and altitude and can be
extended to tables that are functions of four or five independent variables.
